% Contenu partagé (élève/corrigé) — IParcours G2
% !TeX program = lualatex

% Barème indicatif
\bareme{
	\exbadge{1}{5} \quad
	\exbadge{2}{4} \quad
	\exbadge{3}{4} \quad
	\exbadge{4}{3} \quad
	\exbadge{5}{4}
}

% ==================================================================
% PARTIE 1 : À COMPLÉTER DIRECTEMENT SUR LE SUJET
% ==================================================================
\begin{center}
	\Large\textbf{Partie 1 : À compléter directement sur le sujet}
\end{center}
\vspace{2mm}

% =============================
% EXERCICE 1 — Cours et propriétés (5 pts)
% =============================
\begin{exercisebox}{\faIcon{book} Exercice 1 : Cours et propriétés \points{5}}
	\textbf{1.} Énonce le \textbf{théorème de Pythagore}. \points{1}
	
	\anslines[0.4cm]{5}{\begin{minipage}{\linewidth}
			Dans un triangle rectangle, le carré de la longueur de l'hypoténuse est égal à la somme des carrés des longueurs des deux autres côtés.
	\end{minipage}}
	
	\smallspace
	
	\textbf{2.} Complète chaque propriété avec l’égalité de Pythagore correspondante. \points{2}
	\begin{enumerate}[label=\alph*)]
		\item Si un triangle $FQZ$ est rectangle en $Q$ alors\;\onlyeleve{\underline{\hspace{2.2cm}} + \underline{\hspace{2.2cm}} = \underline{\hspace{2.2cm}}} \onlycorrige{$\;FZ^2 = FQ^2 + QZ^2$}.
		\item Si un triangle $JMX$ est rectangle en $X$ alors\;\onlyeleve{\underline{\hspace{2.2cm}} + \underline{\hspace{2.2cm}} = \underline{\hspace{2.2cm}}} \onlycorrige{$\;JM^2 = JX^2 + XM^2$}.
		\item Si un triangle $KNS$ est rectangle en $K$ alors\;\onlyeleve{\underline{\hspace{2.2cm}} + \underline{\hspace{2.2cm}} = \underline{\hspace{2.2cm}}} \onlycorrige{$\;NS^2 = NK^2 + KS^2$}.
		\item Si un triangle $PIO$ est rectangle en $I$ alors\;\onlyeleve{\underline{\hspace{2.2cm}} + \underline{\hspace{2.2cm}} = \underline{\hspace{2.2cm}}} \onlycorrige{$\;PO^2 = PI^2 + IO^2$}.
	\end{enumerate}
	
	\smallspace
	
	\textbf{3.} Énonce la \textbf{réciproque du théorème de Pythagore}. \points{1}
	
	\anslines[0.4cm]{5}{\begin{minipage}{\linewidth}
			Si, dans un triangle, le carré de la longueur du plus grand côté est égal à la somme des carrés des longueurs des deux autres côtés, alors ce triangle est rectangle.
	\end{minipage}}
	
	\smallspace
	
	\textbf{4.} Complète chaque propriété (utiliser la réciproque). \points{1}
	\begin{enumerate}[label=\alph*)]
		\item Si $ST^2 + TJ^2 = SJ^2$ alors le triangle \onlyeleve{\underline{\hspace{1.6cm}}} \onlycorrige{$\;STJ$} est rectangle en \onlyeleve{\underline{\hspace{1cm}}} \onlycorrige{$\;T$}.
		\item Si $MR^2 + MV^2 = RV^2$ alors le triangle \onlyeleve{\underline{\hspace{1.6cm}}} \onlycorrige{$\;MRV$} est rectangle en \onlyeleve{\underline{\hspace{1cm}}} \onlycorrige{$\;M$}.
		\item Si $WX^2 + XY^2 = WY^2$ alors le triangle \onlyeleve{\underline{\hspace{1.6cm}}} \onlycorrige{$\;WXY$} est rectangle en \onlyeleve{\underline{\hspace{1cm}}} \onlycorrige{$\;X$}.
	\end{enumerate}
\end{exercisebox}


\newpage % Sépare clairement les deux parties


% ==================================================================
% PARTIE 2 : À RÉDIGER SUR VOTRE COPIE
% ==================================================================
\begin{center}
	\Large\textbf{Partie 2 : À rédiger sur votre copie}
\end{center}
\vspace{2mm}

% =============================
% EXERCICE 2 — Calcul de longueurs (4 pts)
% =============================
\begin{exercisebox}{\faIcon{calculator} Exercice 2 : Calcul de longueurs \points{4}}
	\textbf{1.} Soit $LMN$ un triangle rectangle en $L$ tel que $LM = \SI{9}{cm}$ et $LN = \SI{12}{cm}$.\\
	Calcule la longueur de l'hypoténuse $[MN]$. \points{2}
	
	\onlycorrige{\vspace{4cm}\par% Espace pour la correction
		Le triangle $LMN$ est rectangle en $L$. D'après le théorème de Pythagore :\\
		$MN^2 = LM^2 + LN^2$\\
		$MN^2 = 9^2 + 12^2$\\
		$MN^2 = 81 + 144 = 225$\\
		$MN = \sqrt{225} = \SI{15}{cm}$.
	}
	
	\smallspace
	
	\textbf{2.} Soit $RST$ un triangle rectangle en $S$ tel que $RT = \SI{10}{cm}$ et $RS = \SI{8}{cm}$.\\
	Calcule la longueur du côté $[ST]$. \points{2}
	
	\onlycorrige{\vspace{4cm}\par% Espace pour la correction
		Le triangle $RST$ est rectangle en $S$. D'après le théorème de Pythagore :\\
		$RT^2 = RS^2 + ST^2$\\
		$10^2 = 8^2 + ST^2$\\
		$100 = 64 + ST^2$\\
		$ST^2 = 100 - 64 = 36$\\
		$ST = \sqrt{36} = \SI{6}{cm}$.
	}
\end{exercisebox}

\exspace

% =============================
% EXERCICE 3 — Vérifier si un triangle est rectangle (4 pts)
% =============================
\begin{exercisebox}{\faIcon{check-double} Exercice 3 : Triangle rectangle ? \points{4}}
	\textbf{1.} Soit $BUT$ un triangle tel que $TU = \SI{5,1}{cm}$, $BU = \SI{7,9}{cm}$ et $BT = \SI{6}{cm}$. Ce triangle est-il rectangle ? Justifie ta réponse. \points{2}
	
	\onlycorrige{\vspace{4cm}\par% Espace pour la correction
		\begin{minipage}{\linewidth}
			Le côté le plus long est $[BU]$. On compare $BU^2$ et $BT^2 + TU^2$ :\\
			D'une part : $BU^2 = 7{,}9^2 = 62{,}41$.\\
			D'autre part : $BT^2 + TU^2 = 6^2 + 5{,}1^2 = 36 + 26{,}01 = 62{,}01$.\\
			Comme $BU^2 \neq BT^2 + TU^2$, l'égalité de Pythagore n'est pas vérifiée. Donc le triangle $BUT$ n’est \textbf{pas} rectangle.
	\end{minipage}}
	
	\smallspace
	
	\textbf{2.} Soit $CAR$ un triangle tel que $RA = \SI{3,3}{cm}$, $AC = \SI{5,6}{cm}$ et $RC = \SI{6,5}{cm}$. Ce triangle est-il rectangle ? Justifie ta réponse. \points{2}
	
	\onlycorrige{\vspace{4cm}\par% Espace pour la correction
		\begin{minipage}{\linewidth}
			Le côté le plus long est $[RC]$. On compare $RC^2$ et $RA^2 + AC^2$ :\\
			D'une part : $RC^2 = 6{,}5^2 = 42{,}25$.\\
			D'autre part : $RA^2 + AC^2 = 3{,}3^2 + 5{,}6^2 = 10{,}89 + 31{,}36 = 42{,}25$.\\
			Comme $RC^2 = RA^2 + AC^2$, d'après la réciproque du théorème de Pythagore, le triangle $CAR$ est rectangle en $A$.
\end{{minipage}}}
\end{exercisebox}

\exspace

% =============================
% EXERCICE 4 — Problème : L'échelle (3 pts)
% =============================
\begin{problembox}{\faIcon{building} Exercice 4 : L'échelle \points{3}}
\begin{minipage}[t]{0.6\linewidth}\vspace{0pt}
	Une échelle de \SI{5}{m} de long est posée contre un mur vertical. Le pied de l'échelle se trouve à \SI{1,4}{m} du pied du mur, sur un sol parfaitement horizontal.
	
	Quelle hauteur atteint le sommet de l'échelle sur le mur ?
\end{minipage} \hfill
\begin{minipage}[t]{0.35\linewidth}\vspace{0pt}
	\centering
	\begin{tikzpicture}[scale=0.7]
		% Mur
		\draw[very thick] (0,0) -- (0,7) node[pos=0.5, left=2mm] {Mur};
		% Sol
		\draw[very thick] (0,0) -- (4,0) node[pos=0.8, below] {Sol};
		% Coordonnées
		\coordinate (sommetMur) at (0,4.8);
		\coordinate (piedEchelle) at (1.4,0);
		\draw[blue, very thick] (piedEchelle) -- (sommetMur); 
		% Annotations
		\draw (0,0) rectangle (0.4,0.4); % Angle droit
		\node[below] at (0.7, 0) {\SI{1,4}{m}};
		\node[above, sloped, blue, xshift=3mm, yshift=2mm] at ( $ (piedEchelle)!0.5!(sommetMur) $ ) {\SI{5}{m}};
		\node[left=2mm] at (0, 2.4) {$h = ?$};
	\end{tikzpicture}
\end{minipage}
\onlycorrige{\vspace{4.5cm}\par% Espace pour la correction
	\begin{minipage}{\linewidth}
		Le mur, le sol et l'échelle forment un triangle rectangle. L'échelle est l'hypoténuse. Soit $h$ la hauteur atteinte.\\
		D'après le théorème de Pythagore :\\
		$5^2 = h^2 + 1{,}4^2$\\
		$25 = h^2 + 1{,}96$\\
		$h^2 = 25 - 1{,}96 = 23{,}04$\\
		$h = \sqrt{23{,}04} = \SI{4,8}{m}$.\\
		Le sommet de l'échelle atteint une hauteur de \SI{4,8}{m}.
\end{minipage}}
\end{problembox}

\exspace

% =============================
% EXERCICE 5 — Problème : Table de chevet (4 pts)
% =============================
\begin{problembox}{\faIcon{cubes} Exercice 5 : La table de chevet \points{4}}
	% --- Présentation du problème (Texte à gauche, Figure à droite) ---
	\begin{minipage}[t]{0.65\linewidth}\vspace{0pt}
		Un designer a imaginé une table de chevet. On modélise sa structure métallique ci-contre. Les segments \([JL]\) et \([KI]\) sont \textbf{perpendiculaires} en \(M\).
		
		On donne : $JM = \SI{38,4}{cm}$, $MK = \SI{28,8}{cm}$ et $MI = \SI{51,2}{cm}$.\\
		
		\vspace{2mm}
			% --- Questions sur toute la largeur ---
		\textbf{1.} Calcule les longueurs des montants $[JK]$ et $[JI]$. \points{2}\\
		
			\textbf{2.} Le plateau, représenté par le segment $[LK]$, doit être perpendiculaire au montant $[JK]$. Est-ce le cas ? Justifie. \points{2}
		\onlycorrige{\vspace{4cm}\par% Espace pour la correction
			On doit vérifier si le triangle $JKL$ est rectangle en $J$. On a $JK=48$. Il manque $JL$ et $LK$.
			En supposant que la question visait à savoir si le plateau $[IK]$ est perpendiculaire à la base $[JL]$, on vérifie si $IJK$ est rectangle en $J$.
			On a $KI = KM + MI = 28{,}8 + 51{,}2 = \SI{80}{cm}.$\\
			Dans le triangle $JKI$, le plus long côté est $[KI]$.
			D'une part : $KI^2 = 80^2 = 6400$.\\
			D'autre part : $JK^2 + JI^2 = 48^2 + 64^2 = 2304 + 4096 = 6400$.\\
			Comme $KI^2 = JK^2 + JI^2$, d'après la réciproque, le triangle $JKI$ est rectangle en $J$.}
	\end{minipage}
	\hfill % Espace élastique qui pousse la figure vers la droite
	\begin{minipage}[t]{0.3\linewidth}\vspace{0pt}\centering
		\def\chevetimg{../images/fig_table_chevet_controle_pythagore}
		\IfFileExists{\chevetimg.pdf}{\includegraphics[width=\linewidth]{\chevetimg.pdf}}{%
			\IfFileExists{\chevetimg.png}{\includegraphics[width=\linewidth]{\chevetimg.png}}{%
				\IfFileExists{\chevetimg.jpg}{\includegraphics[width=\linewidth]{\chevetimg.jpg}}{%
					\begin{tikzpicture}[x=1mm,y=1mm]
						\draw[gray] (0,0) rectangle (70,45);
						\node at (35,22.5) {\small Image manquante};
					\end{tikzpicture}%
				}%
			}%
		}
	\end{minipage}
	
	%\vspace{4mm} % Petit espace vertical avant les questions
	

	\onlycorrige{\vspace{2.5cm}\par% Espace pour la correction
		Dans le triangle $JKM$ rectangle en $M$ : $JK^2 = JM^2+MK^2 = 38{,}4^2+28{,}8^2 = 1474.56 + 829.44 = 2304$. Donc $JK = \sqrt{2304} = \SI{48}{cm}$.\\
		Dans le triangle $JIM$ rectangle en $M$ : $JI^2 = JM^2+MI^2 = 38{,}4^2+51{,}2^2 = 1474.56 + 2621.44 = 4096$. Donc $JI = \sqrt{4096} = \SI{64}{cm}$.
	}
\end{problembox}